\chapter{Commenting out Code}
Sometimes you want to \bsq{comment out} code that should (no longer) be compiled and/or executed, instead of removing it completely.

First, your production code should not contain any code that is commented out. Nevertheless, during development it is sometimes helpful that you have this option.

Use the \bdq{\lstinline|//|} delimiter to comment out a complete line, only a partial line, or a whole bunch of consecutive lines of code.

Example:
\begin{lstlisting}
if( foo > 1 )
{
    //---* Do a double-flip *----------------------------------------
    …
}
else
{
    return false;  // Explain why here.
}
//if( bar > 1 )
//{
//    //---* Do a triple-flip *--------------------------------------
//    …
//}
//else
//{
//    return false;
//}
\end{lstlisting}

It is not a good idea to use a block comment (\bdq{\lstinline|/*…*/|}) for commenting out sections of code, because you cannot have a block comment inside a block comment.

The easiest way to comment out a selected code block this way in the Eclipse IDE\index{Eclipse!IDE} is to use \bdq{\texttt{Source|\allowbreak Toggle Comment}} from the menu or the short key \bdq{\texttt{CTRL+/}}\footnote{On a German keyboard, it is \bdq{7} instead of \bdq{/}.}.

In IntelliJ~IDEA\index{IntelliJ IDEA}, it is the menu command \bdq{\texttt{Code|\allowbreak Comment with Line Comment}} or the short key \bdq{\texttt{CTRL+/}}\footnote{On a German keyboard, it is \bdq{÷} on the numeric keypad instead of \bdq{/}.}.

The configuration setting \bdq{\texttt{File|\allowbreak Settings…|\allowbreak Editor>\allowbreak Inspections>\allowbreak Code Maturity>\allowbreak Commented out code}} in IntelliJ~IDEA\index{IntelliJ IDEA} allows that a warning or an error is generated for code that was commented out.

%==============================================================================
% End of file!!
%==============================================================================

